\section{Потребителски сценарии и места за screenshots}

Една от важните функции на техническата дипломна документация е да направи връзка между архитектурното описание и реалната употреба на системата. Поради тази причина в настоящата глава са представени основните потребителски сценарии в последователен демонстрационен вид. Всеки сценарий е придружен от описание на очакваното поведение на системата и от предварително предвидено място за screenshot от работещото приложение.

Този подход има две предимства. От една страна, текстът действа като кратко ръководство за демонстрация. От друга страна, структурата позволява реални екранни снимки да бъдат добавени на по-късен етап без нужда от пренаписване на документа. Така дипломната работа остава едновременно съдържателно завършена и practically подготвена за визуално допълване.

\subsection{Сценарий 1: Начална страница и избор на език}

Демонстрацията започва от началната страница на приложението. Тя има ролята на публична входна точка и предоставя достъп до регистрация, вход и избор на език. При липса на активна сесия системата не позволява използване на вътрешните функции, а се държи като портал към удостоверителния процес. Визуално началната страница съдържа хедър, навигационна лента, hero секция и основни бутони за навигация.

От потребителска гледна точка това е правилен първи екран, защото поставя акцент върху най-важните начални действия: регистриране или вход. Паралелно с това наличието на връзки за \texttt{EN} и \texttt{BG} в навигацията демонстрира, че приложението поддържа базова локализация още от първия момент на взаимодействие.

\screenshotplaceholder{Начална страница и езиков превключвател}{Тук да се постави screenshot на началната страница, на който ясно да се виждат навигационната лента, hero секцията и връзките за смяна на езика.}

\subsection{Сценарий 2: Регистрация на нов потребител}

Следващата естествена стъпка е регистрацията. Потребителят отваря формата за регистрация и попълва име, фамилия, имейл и парола. При коректно подадени данни системата създава нов профил и инициира изпращане на потвърждаващ e-mail. Ако входът не е коректен или ако вече съществува потребител със същия имейл, се визуализира подходяща обратна връзка.

Този сценарий има ключова демонстрационна стойност, защото показва едновременно работа с форма, server-side обработка, persistence операция и интеграция с външен mail компонент. Именно тук се вижда, че проектът надхвърля чисто визуален прототип и реализира реален поток по създаване на акаунт.

\screenshotplaceholder{Форма за регистрация}{Тук да се постави screenshot на екрана за регистрация с попълнени примерни данни или с демонстрация на валидационни съобщения при невалиден вход.}

\subsection{Сценарий 3: Потвърждение на акаунт и вход}

След изпращане на регистрационната форма потребителят следва да потвърди своя акаунт чрез връзката, получена по електронна поща. В практическа демонстрация може да се покаже както самият имейл, така и резултатът от активиране на линка -- преминаване към форма за вход. След потвърждение потребителят въвежда имейл и парола, а системата проверява валидността на данните и създава активна сесия.

От гледна точка на потребителската история това е преходът от публичната към вътрешната част на приложението. От този момент нататък системата започва да показва допълнителни навигационни елементи, които са достъпни само при активна сесия.

\screenshotplaceholder{Вход в системата след потвърдена регистрация}{Тук да се постави screenshot на login формата или на резултата след успешно удостоверяване, например пренасочване към каталога.}

\subsection{Сценарий 4: Профил на потребителя и списък с фирми}

След вход един от най-важните екрани е профилът. Той позволява на потребителя да прегледа своите основни данни и при необходимост да ги редактира. Допълнително в страничната част на екрана се визуализира списък от свързани фирми и е наличен бутон за добавяне на нова.

Този екран е показателен за комбиниран характер на интерфейса: той не е само форма за редакция, а едновременно работна зона, която обединява идентичност на потребителя и свързана бизнес информация. При демонстрация е подходящо да се покаже както нормалното състояние на профила, така и пример за съобщение при неправилно въведена стара парола или несъвпадение между новата и потвърдената парола.

\screenshotplaceholder{Профил на потребителя}{Тук да се постави screenshot на профилния екран с визуализирани полета за име, имейл, парола и списък с фирми на потребителя.}

\subsection{Сценарий 5: Добавяне на фирма}

Формата за добавяне на фирма е логически подготвителна стъпка за последващото въвеждане на продукти. Потребителят посочва име на фирмата, идентификационни данни, адрес и тип на регистрация. След успешен запис фирмата става достъпна в списъка от фирми и може да бъде избирана от формата за нов продукт.

Този сценарий е особено полезен за демонстрация пред комисия, защото показва, че приложението не работи с изолирани продукти, а поддържа по-реалистична домейн структура, в която продуктът се асоциира с фирма. По този начин каталогът се обвързва с източник на данни, а не с анонимни артикули.

\screenshotplaceholder{Форма за добавяне на фирма}{Тук да се постави screenshot на екрана за добавяне на фирма, включително полетата за име, MOL, VAT, адрес и статус на регистрация.}

\subsection{Сценарий 6: Добавяне на продукт}

След като съществува поне една фирма, потребителят може да премине към добавяне на продукт. Екранът за това включва полета за име, URL към изображение, цена, количество, категория и избор на фирма. Категорията се избира от предварително дефиниран enum списък, а фирмата -- от асоциираните с потребителя компании.

Този сценарий показва няколко важни характеристики на приложението. Първо, демонстрира комбиниране на текстови, числови и избираеми полета. Второ, показва как входните данни се валидират още на ниво форма. Трето, илюстрира как домейн връзките вече се материализират в UI: потребителят не въвежда фирма на свободен текст, а я избира от наличен набор.

\screenshotplaceholder{Добавяне на нов продукт}{Тук да се постави screenshot на формата за добавяне на продукт, на който ясно да се виждат изборът на категория, фирма и валидационните полета.}

\subsection{Сценарий 7: Каталог и филтриране по категория и цена}

След добавяне на продукти системата може да бъде демонстрирана чрез каталога. Това е един от най-визуалните екрани в приложението. Потребителят вижда странично меню с категории и зона с продуктови карти. Допълнително JavaScript логиката използва REST крайната точка за максимална цена, за да конфигурира ценови плъзгач за клиентско филтриране.

При демонстрация е подходящо да се покажат поне три действия: отваряне на общия каталог, преминаване към конкретна категория и стесняване на видимите продукти чрез ценовия диапазон. Този сценарий представя най-добре съчетанието между сървърно рендиран интерфейс и ограничена клиентска динамика.

\screenshotplaceholder{Каталог с продукти}{Тук да се постави screenshot на каталога с няколко продукта, страничното меню с категории и активния ценови филтър.}

\subsection{Сценарий 8: Детайли на продукт и добавяне в количка}

След избор на конкретен артикул потребителят преминава към страницата с детайли. Там се визуализират изображение, наименование, цена и текуща наличност. Потребителят въвежда желаното количество и го добавя в количката. Ако количеството е невалидно или надвишава наличността, системата връща съобщение за грешка.

Този екран има съществена демонстрационна роля, защото именно тук се вижда преходът от пасивно разглеждане към активна транзакционна интеракция. От UX гледна точка формата е минималистична и фокусира вниманието върху единичния продукт.

\screenshotplaceholder{Детайлен изглед на продукт}{Тук да се постави screenshot на страницата с детайли на продукт, включително полето за количество и бутона за добавяне в количката.}

\subsection{Сценарий 9: Количка и финализиране на поръчка}

След добавяне на продукти в количката потребителят отваря съответния екран, където вижда избраните артикули, броя им, цената на всеки запис, общата стойност, таксата за доставка и крайна сума. В summary панела се попълва адресът, след което се извършва checkout.

Този сценарий е особено важен, защото демонстрира, че приложението не спира до каталог и визуализация, а реализира реален край на потребителската цел. При защита на дипломна работа именно тази страница често е сред най-силните екрани за демонстрация, защото интегрира back-end изчисления, визуално обобщение и form-based завършване на операцията.

\screenshotplaceholder{Количка и checkout екран}{Тук да се постави screenshot на количката, така че да се виждат продуктите, общата сума, таксата за доставка и полето за адрес.}

\subsection{Сценарий 10: Екран за успешна поръчка и история}

След завършване на покупката потребителят се пренасочва към потвърждаващ екран, а впоследствие може да отвори своята история на поръчките. Историческият изглед показва списък с реализираните покупки, включително идентификатор, имейл, адрес, дата и цена. Наличен е и бутон за изчистване на историята.

Този етап е важен за демонстрацията, защото визуализира персистентния резултат от целия потребителски процес. Докато каталожните и checkout екраните показват текущо състояние, историята показва, че системата поддържа и времева следа за вече реализирани действия.

\screenshotplaceholder{История на поръчките}{Тук да се постави screenshot на екрана с история на поръчките, включително таблицата с дата, адрес и цена на направените покупки.}

\subsection{Сценарий 11: Демонстрация на локализацията}

Допълнителен, но важен демонстрационен сценарий е смяната на езика. Тъй като системата поддържа базова английска и българска локализация, е полезно в дипломната работа да се покаже поне един и същи екран в двата режима. Това дава визуално доказателство, че локализационната конфигурация не е само декларативно налична, а реално участва в поведението на приложението.

При добра демонстрация могат да бъдат сравнени например началната страница, формата за добавяне на продукт или екранът за профил в двата езикови режима. Подобна съпоставка е особено ценна, защото показва и ограниченията на текущата локализация -- къде тя е пълна и къде остават твърдо кодирани текстове.

\screenshotplaceholder{Съпоставка между английски и български интерфейс}{Тук да се постави двойка screenshots или комбиниран кадър, показващ един и същ екран на английски и на български език.}

\subsection{Препоръки за създаване на screenshots}

За да бъде визуалната част на дипломната работа достатъчно качествена, е препоръчително screenshots да бъдат направени при следните условия:

\begin{itemize}[leftmargin=*]
    \item браузърът да бъде отворен на достатъчно голяма резолюция, така че всички основни елементи да се виждат ясно;
    \item демонстрационните данни да бъдат подготвени предварително, за да няма празни списъци или непълни форми;
    \item при нужда да се подберат данни с по-описателни имена на фирми и продукти;
    \item да се избягва показване на чувствителни конфигурационни данни, локални пароли и реални секрети;
    \item да се поддържа последователен визуален стил между всички използвани изображения.
\end{itemize}

Тези препоръки са важни, защото screenshots не са просто украшение към документа. Те са доказателствен материал, че описаните сценарии са реално демонстрируеми и че системата има работещ потребителски интерфейс.

\subsection{Обобщение}

Настоящата глава свързва архитектурното и имплементационното описание на проекта с неговото реално потребителско поведение. Представените сценарии обхващат целия жизнен цикъл на работа със системата -- от началния достъп и регистрацията, през управление на фирми и продукти, до каталога, количката, checkout процеса и историята на поръчките. Оставените места за screenshots осигуряват готова рамка за визуално завършване на дипломната работа.

По този начин документът не остава само аналитичен, а се превръща и в практическо ръководство за демонстрация на проекта. Следващата глава преминава към по-формална верификация и оценка на полученото решение.
